\chapter{Scope, Motivation and Project Status / Escopo, Motivação e Estado do Projeto}

\section*{English}
The RADOME project proposes a distributed passive electromagnetic sensing network based on geodesic and conformal radomes installed at elevated sites. Each station combines multiband receiving apertures, local digitization, event-oriented processing, precise timing, calibration and optical networking. Several stations cooperate to estimate angle of arrival (AOA), time difference of arrival (TDOA), frequency difference of arrival (FDOA), Doppler and polarization.

The network may exploit known opportunistic illuminators, including cellular infrastructure, broadcast transmitters, satellites and other authorized RF sources. In bistatic operation, the direct reference path is compared with the path involving a target and a receiver. The system is conceptual: performance, coverage, locations and accuracy must be demonstrated experimentally and are not operational guarantees.

The project objectives are to develop a measurable multiband architecture, validate a three-node VHF/UHF demonstrator, characterize end-to-end timing and RF delays, preserve event windows locally, and integrate the electromagnetic platform with resilient energy, thermal, communications and site infrastructure.

This article, assembled from \texttt{projetov1.tex}, is the authoritative technical document. ``Document version 1.1'' identifies the controlled publication, while ``architecture revision 3'' identifies the current conceptual system architecture; these identifiers are independent. Earlier Markdown, PDF and standalone LaTeX documents are historical inputs unless a current chapter explicitly incorporates their content. Quantitative values are controlled in \texttt{PARAMETERS.md}, with status proposed, simulated, measured or derived, and architecture decisions are recorded in \texttt{DECISIONS.md}. In case of conflict, an approved parameter-register entry and this article take precedence over historical material.

\section*{Português}
O projeto RADOME propõe uma rede distribuída de sensoriamento eletromagnético passivo baseada em radomes geodésicos e conformais instalados em sítios elevados. Cada estação combina aberturas receptoras multifaixa, digitalização local, processamento orientado a eventos, referência temporal, calibração e rede óptica. Várias estações cooperam para estimar ângulo de chegada (AOA), diferença de tempo de chegada (TDOA), diferença de frequência de chegada (FDOA), Doppler e polarização.

A rede pode explorar iluminadores de oportunidade conhecidos, incluindo infraestrutura celular, radiodifusoras, satélites e outras fontes de RF autorizadas. Na operação biestática, o caminho direto de referência é comparado ao caminho que envolve o alvo e o receptor. O projeto é conceitual: desempenho, cobertura, locais e precisão devem ser demonstrados experimentalmente e não constituem garantias operacionais.

Os objetivos são desenvolver uma arquitetura multifaixa mensurável, validar um demonstrador de três nós VHF/UHF, caracterizar os atrasos temporais e RF ponta a ponta, preservar janelas de eventos localmente e integrar a plataforma eletromagnética à infraestrutura energética, térmica, de comunicações e de sítio.

Este artigo, montado a partir de \texttt{projetov1.tex}, é o documento técnico autoritativo. A ``versão do documento 1.1'' identifica a publicação controlada, enquanto a ``revisão da arquitetura 3'' identifica a arquitetura conceitual vigente; os dois identificadores são independentes. Documentos Markdown, PDF e LaTeX autônomos anteriores são fontes históricas, salvo quando seu conteúdo for incorporado explicitamente por um capítulo vigente. Os valores quantitativos são controlados em \texttt{PARAMETERS.md}, com estado proposto, simulado, medido ou derivado, e as decisões de arquitetura são registradas em \texttt{DECISIONS.md}. Em caso de conflito, uma entrada aprovada do registro de parâmetros e este artigo prevalecem sobre o material histórico.
